\section{WebStorm}
WebStorm jest jednym z najbardziej popularnych jak i najbardziej zaawansowanych środowisk programistycznych (IDE). Głównie służy do obsługi projektów webowych.
\subsection{Działanie i zalety WebStorm}
\begin{enumerate}
    \item \textbf{Wsparcie dla wielu językow programowania jak i narzędzi:}
    WebStorm zapewnia wsparcie dla języków programowania takich jak JavaScript, TypeScript jak i Node.js. Obsługuje również frameworki i biblioteki takie jak React. Dodatkowo obsługuje HTML, CSS i inne narzędzia umożliwiające budowanie aplikacji webowych.
    \item \textbf{Inteligentne podpowiedzi:} inteligentne mechanizmy analizy kodu jak i wbudowane modele AI dostarczają podpowiedzi i umożliwiają szybsze pisanie kodu. Przykładem są podpowiedzi zmiennych, importów jak i podstawowych funkcji np. handlerów od eventów.
    \item \textbf{Integracja z systemem kontroli wersji Git:} WebStorm posiada wrapper odpowiedzialny za integrację z Gitem, umożliwia ona przede wszystkim:
    \begin{itemize}
        \item Tworzenie i zarządzanie gałęziami,
        \item Obsługę commitów, w tym wysyłania jak i pobierania ich z serwera,
        \item Rozwiązywanie konfliktów przy merge,
        \item Przeglądanie historii zmian.
    \end{itemize}
\end{enumerate}
